% Reciprocal proofs of the declared modern reconstructions; no philosophical admission.
\begin{annexitem}{annex:p15.popper.package}{thm:p15-popper-package}{The tested conjunction and the countervaluations}
\begin{proof}
Let a classical valuation satisfy \(\chi,\neg E\).
If it also satisfied \(H,A,I\), the antecedent of \(\chi\)
would be true and the consequent false. This contradicts
satisfaction of the conditional. The conjunction is therefore false.
With \(A,I\) additionally true, \(H\) must be false.

For the nonentailments set \(E\) false. Each of
\[
(H,A,I)=(\mathsf{true},\mathsf{false},\mathsf{true}),\quad
(\mathsf{false},\mathsf{true},\mathsf{true}),\quad
(\mathsf{true},\mathsf{true},\mathsf{false})
\]
makes the conditional true by making its antecedent false.
The first is a countervaluation to entailment of \(\neg H\);
the second to entailment of \(\neg A\); the first or second
to entailment of \(\neg I\).
The third explicitly gives initial-condition failure while the
other two components remain true.
For a positive test set \(H\) false and \(A,I,E\) true.
All the displayed positive premises hold, but \(H\) does not.

Substituted classical formulas obey the same Boolean truth clauses
under every interpretation, so the entailments persist.
Independent assignability is needed only to realize the stated
countervaluations; no nonentailment for every substitution was claimed.
\end{proof}
\end{annexitem}

\begin{annexitem}{annex:p15.popper.thread}{thm:p15-popper-thread}{Prediction and three complete diagnostic models}
\begin{proof}
In a model of \(A,I\), calibration gives
\(L(c)=r_L(c)=2\) and \(S(c)=r_S(c)=1\).
Since \(2>1\), the law \(H\) gives \(\mathsf{Br}(c)\).
Detector equivalence then gives \(\mathsf R(c)\), namely \(E\).

For all three models let \(U=\{c\}\), interpret the constant
by its only element, and fix
\(r_L(c)=2,r_S(c)=1,\mathsf R(c)\) false.
All functions below have positive values and these choices satisfy
\(I,\neg E\).

For law failure set \(L(c)=2,S(c)=1,\mathsf{Br}(c)\) false.
Both calibration equalities hold. Both sides of the detector
equivalence are false, so it holds and \(A\) is true.
The law's antecedent is true and its consequent false, so \(H\)
is false.

For strength-report failure set \(L(c)=2,S(c)=3\) and
\(\mathsf{Br}(c)\) false. The law's only implication has
false antecedent \(2>3\), so \(H\) is true.
Load calibration and detector equivalence hold.
Strength calibration fails because \(r_S(c)=1\neq3=S(c)\).

For detector failure set \(L(c)=2,S(c)=1\) and
\(\mathsf{Br}(c)\) true. The law and both calibration equalities
hold. Detector equivalence fails because actual breaking is true
and reported breaking false.
The complete displayed observation record is identical in all three.
It therefore cannot uniquely entail any of the three different
diagnoses.
\end{proof}
\end{annexitem}

\begin{annexitem}{annex:p15.popper.finite-record}{thm:p15-popper-finite-record}{Unnamed witnesses and exhaustively enumerated domains}
\begin{proof}
For \(n\geq0\) use domain \(\{0,1,\ldots,n\}\), with \(a_i\)
denoting \(i\) for \(1\leq i\leq n\).
In both structures let \(\mathsf{Rv}\) hold everywhere.
In the first let \(\mathsf{Wh}\) hold everywhere; in the second
let it hold exactly on \(\{1,\ldots,n\}\).
The constants are distinct and every named object is a white raven
in both structures. All indicated predicate facts about those
objects agree.
The first structure satisfies \(H_{\rm rw}\).
The second refutes it at the unnamed element \(0\).
For \(n=0\) the record is empty but both domains remain nonempty,
so the construction still applies.

In any model of
\(\mathsf{Rv}(b)\land\neg\mathsf{Wh}(b)\), the denotation of \(b\)
violates the universal implication. Thus \(\neg H_{\rm rw}\)
holds. This proves a conditional consequence, not an empirical
identification of the object.

For finite negative \(P\)-observations use the same domain and
constants. Give \(P\) empty extension in one structure and
extension \(\{0\}\) in another. Both satisfy every \(\neg P(a_i)\),
but only the second satisfies \(\exists xP(x)\).
Hence the negative record does not entail nonexistence.
If \(P(b)\) holds, the denotation of \(b\) directly witnesses
\(\exists xP(x)\).

Finally assume \(n\geq1\) and the exhaustive-domain premise.
Universal instantiation gives every displayed implication at \(a_i\)
from \(H_{\rm rw}\).
Conversely, suppose all those implications hold and let \(u\)
be any domain element. Exhaustiveness gives \(u=a_i\) for some
\(i\). If \(\mathsf{Rv}(u)\) holds, equality of denotations and
that instance give \(\mathsf{Wh}(u)\).
Since \(u\) was arbitrary, \(H_{\rm rw}\) holds.
This last equivalence does not require unique names.
\end{proof}
\end{annexitem}

\begin{annexitem}{annex:p15.popper.falsifier-order}{thm:p15-popper-falsifier-order}{Inclusion, equality despite strengthening and incomparable tests}
\begin{proof}
Take \(b\in F_{\mathcal K,\mathcal B}(T_2)\).
If some \(\mathcal M\in\mathcal K\) satisfied \(T_1\land b\),
the consequence hypothesis would make it satisfy \(T_2\land b\),
contrary to the choice of \(b\).
Therefore \(b\in F_{\mathcal K,\mathcal B}(T_1)\).
The specified strictness witness belongs to this latter class
but not the former, so the inclusion is proper.

For equality despite strict strengthening, take all Boolean
valuations of \(p,q,r\) as \(\mathcal K\), and set
\[
\mathcal B=\{\neg p,r,\neg r\},\qquad
T_1=p\land q,\quad T_2=p.
\]
Every admitted test is individually satisfiable.
Both theories contradict \(\neg p\).
Both are compatible with \(r\), by setting \(p,q,r\) true,
and with \(\neg r\), by setting \(p,q\) true and \(r\) false.
Both falsifier classes are exactly \(\{\neg p\}\).
Yet \(T_1\) entails \(T_2\), and \(p\) true, \(q\) false
refutes the converse. This explicitly restricted observational
vocabulary does not test the extra \(q\).

A tautology is compatible with each \(b\), since the definition
requires a \(\mathcal K\)-model of each \(b\).
For example, \(E\lor\neg E\) is true under every Boolean valuation:
either \(E\) is true or its negation is true.
A contradiction is compatible with none.
The individual satisfiability assumption is essential:
an impossible test statement is not an informative falsifier
of a tautology.

For incomparability take all Boolean valuations of \(p,q\) and
\(\mathcal B=\{p,\neg p,q,\neg q\}\).
The falsifier class of \(p\) is exactly \(\{\neg p\}\),
and that of \(q\) exactly \(\{\neg q\}\):
each other candidate has an explicitly compatible valuation.
These distinct singleton classes neither include one another.
No total ordering follows from inclusion alone.
\end{proof}
\end{annexitem}

\begin{annexitem}{annex:p15.popper.probability-order}{thm:p15-popper-probability-order}{Conditional inclusion under a stipulated probability law}
\begin{proof}
Intersections and differences of the stipulated events are measurable.
The disjoint sets \(V_1\cap D\) and
\((V_2\setminus V_1)\cap D\) have union \(V_2\cap D\).
Finite additivity gives
\[
\mathbb P(V_2\cap D)
=\mathbb P(V_1\cap D)+
 \mathbb P((V_2\setminus V_1)\cap D).
\]
Divide by the positive number \(\mathbb P(D)\) and use
event conditioning to obtain the asserted difference.
The right side is nonnegative, proving the inequality.
It is positive exactly when the displayed difference event
has positive probability, not merely when inclusion is proper.

For equality despite proper inclusion let
\(\Omega=\{a,b\}\) have masses \(1/2,1/2\).
Take \(V_1=\{a\},V_2=\Omega,D=\{a\}\).
Both conditional probabilities equal one.
For \(\mathbb P(D)=0\), the ratio definition used here is
unavailable; no assertion for that boundary was made.
\end{proof}
\end{annexitem}

\begin{annexitem}{annex:p15.popper.coin}{thm:p15-popper-coin}{A positive-probability realization of the rare observation}
\begin{proof}
There are \(2^n\) binary length-\(n\) sequences.
Mass \(2^{-n}\) on each gives a law because
\(2^n2^{-n}=1\). In particular the all-one sequence is
one of its points, with positive singleton mass \(2^{-n}\).
Taking that point as the realized outcome leaves the law unchanged
and realizes the observation under the law.
This is a consistency witness for law and occurrence.

At \(n=10\), \(2^{10}=1024\), and
\(1/1024<1/20=0.05\) because \(20<1024\).
A rejection convention is a declared map from observed outcomes
to decisions. It does not remove the selected point from the
law's support. Thus that decision and calculation do not entail
logical incompatibility of the uniform law with the observation.
Nothing in the construction establishes that an actual
experimental mechanism has this law.
\end{proof}
\end{annexitem}
